\documentclass{ntnuthesis}
\usepackage[utf8]{inputenc}
\usepackage[T1]{fontenc}
\usepackage{geometry}
\usepackage{graphicx}
\usepackage{float}
\usepackage{hyperref}
\usepackage[english,indonesian]{babel}

\title{Optimasi Pemrosesan Template LaTeX Menggunakan Pendekatan Berbasis AI}
\author{Academic Support Agent}
\date{\today}

\begin{document}
\maketitle

\begin{abstract}

\end{abstract}

\chapter{PENDAHULUAN}


\chapter{TINJAUAN PUSTAKA}


% Original content (commented out):
% \documentclass[british,titlepage]{ntnuthesis}

\title{An NTNU Thesis \LaTeX{} Document Class}
\shorttitle{An NTNU Thesis Document Class}
\author{Community of Practice in Computer Science Education at NTNU}
\shortauthor{CoPCSE$@$NTNU}
\date{CC-BY \ntnuthesisdate}

\addbibresource{thesis.bib}

\input{glossary.tex} % add glossary and acronym lists before document

\begin{document}

\input{chapters/0a-abstract.tex}
\input{chapters/0b-sammendrag.tex}

\tableofcontents
\listoffigures
\listoftables
\lstlistoflistings

\printglossary[type=\acronymtype] % Print acronyms
\printglossary                    % Print glossary

\input{chapters/1-introduction}
\input{chapters/2-usage.tex}
\input{chapters/3-structure.tex}
\input{chapters/4-conclusion.tex}

\chapter*{\bibname}
\printbibliography[heading=none]

\input{chapters/papers.tex}

\appendix
\input{appendices/a-appendix.tex}

\end{document}


\end{document}